\chapter{From genomic representations to informative markers}
\label{chap:feature_importance}

Learning a compact representation answers only part of the problem. A
latent vector may support breed classification while still leaving
open the question of which original SNP markers carried the relevant
information.

This chapter describes how the thesis moves from representation
learning back to the genomic input space. The central idea is to
construct global SNP rankings using several different notions of
relevance and then judge those rankings through their downstream
utility. A marker receives value not merely because an attribution
algorithm assigns it a large score, but because a panel built from
highly ranked markers should preserve useful breed-discriminative
information.

\section{Why marker importance is a separate problem}
\label{sec:why_marker_importance}

Prediction and interpretation are related but different goals. A
classifier can achieve good performance by using many correlated
features without making it obvious which of them are individually
important. In genomic data this issue is amplified by LD: neighboring
markers can contain overlapping information, so multiple plausible
rankings may exist.

The thesis therefore avoids treating any single importance method as
a biological ground truth. Instead, different ranking strategies are
viewed as hypotheses about which parts of the genotype matrix are most
useful for population discrimination.

These hypotheses are tested by constructing reduced panels and
re-evaluating the downstream classification task.

\begin{figure}[htbp]
\centering
\begin{tikzpicture}[
    >=Latex,
    font=\small,
    box/.style={draw, semithick, rounded corners=2pt, align=center,
                minimum height=12mm, minimum width=31mm, inner sep=4pt, fill=black!2},
    key/.style={draw, semithick, rounded corners=2pt, align=center,
                minimum height=10mm, minimum width=62mm, inner sep=4pt, fill=black!7},
    arr/.style={->, semithick}
]
\node[key] (x) at (0,3.05) {Training genotype matrix ($N \times P$)};
\node[box] (fst) at (-5.25,1.35) {$F_{ST}$\\population differentiation};
\node[box] (rf)  at (-1.75,1.35) {Random Forest\\Gini importance};
\node[box] (shap) at (1.75,1.35) {SHAP\\VMGP breed logits};
\node[box] (ig)  at (5.25,1.35) {Integrated Gradients\\contrastive centroid};
\node[key] (rank) at (0,-0.55) {Global SNP ranking $s_1 \ge \cdots \ge s_P$};
\node[key] (panels) at (0,-2.15) {Nested panels $P_K$, $K \in \{50,\ldots,5000\}$};
\node[key] (eval) at (0,-3.75) {Downstream Random Forest, validation-only evaluation};

\draw[arr] (x) -- (fst);
\draw[arr] (x) -- (rf);
\draw[arr] (x) -- (shap);
\draw[arr] (x) -- (ig);
\draw[arr] (fst) -- (rank);
\draw[arr] (rf) -- (rank);
\draw[arr] (shap) -- (rank);
\draw[arr] (ig) -- (rank);
\draw[arr] (rank) -- (panels);
\draw[arr] (panels) -- (eval);
\end{tikzpicture}

\caption{Marker-ranking families compared in RQ4. Each family defines
marker relevance through a different principle: population
differentiation, predictive impurity importance, additive model
attribution, and gradient-based attribution. All rankings are
constructed inside the training partition and are converted into the
same nested panel sizes, whose utility is judged by the same downstream
classifier \cite{weir1984estimating,breiman2001random,
lundberg2017unified,sundararajan2017axiomatic}.}
\label{fig:ranking_families}
\end{figure}


Figure~\ref{fig:ranking_families} summarizes the four ranking families
that are constructed and compared in the experimental evaluation.

\section{From local attributions to a global SNP ranking}
\label{sec:local_to_global_ranking}

Neural attribution methods usually explain a particular model output
for a particular individual. Marker reduction, however, requires a
single ranking that can be applied to the dataset.

Let $a_{i,j}$ denote the attribution assigned to SNP $j$ for
individual $i$ under a specified target. A generic global relevance
score can be written as

\begin{equation}
    s_j
    =
    \mathcal{A}
    \left(
        a_{1,j},
        a_{2,j},
        \ldots,
        a_{N,j}
    \right),
\end{equation}

where $\mathcal{A}$ is an aggregation rule. Markers are then sorted
according to $s_j$ and nested panels are constructed from the
top-ranked loci.

The exact aggregation rule matters. A signed mean, an absolute mean,
a class-balanced aggregation, and a maximum attribution can produce
different interpretations. For this reason, aggregation is treated as
part of the marker-ranking method and must be fixed before definitive
RQ4 evaluation.

All neural rankings follow fixed, executable aggregation rules. For the
SHAP-based ranking of the variational model, the global score of a
marker is the mean absolute attribution over the evaluation subset and
over all breed outputs,
\begin{equation}
    s_j
    =
    \frac{1}{C\,N_{\mathrm{eval}}}
    \sum_{c=1}^{C}
    \sum_{i=1}^{N_{\mathrm{eval}}}
    \left|
        \phi_{i,j,c}
    \right| .
\end{equation}
For the contrastive model, Integrated Gradients values are reduced to a
per-marker value per sample as the $L_1$ sum over the four genotype
channels; the per-sample values are then averaged within each breed and
finally averaged with equal weight across breeds. Both aggregation
rules are implemented in the shared ranking module and are not modified
after inspecting downstream panel performance.

\section{Explaining the variational breed predictor}
\label{sec:vmgp_marker_attribution}

The VMGP-inspired model contains a supervised breed-prediction head.
This gives the attribution problem a direct target: changes in the SNP
input can be related to the output associated with a particular breed
prediction.

SHAP \cite{shapley1953value, lundberg2017unified} is considered as a
way to decompose model predictions into feature-level contributions.
For a fixed model and a defined background distribution, SHAP assigns
each input feature a contribution relative to the reference
prediction.

In the genomic setting, the resulting values are first associated with
individual SNPs and individuals. They must then be aggregated across
observations and, where appropriate, across breed-specific outputs
before a global marker ranking is obtained.

The important methodological point is that the ranking must be
constructed using only the training information available inside the
current experimental partition. The held-out test set cannot be used as
the SHAP background, for attribution aggregation, or for deciding
which markers should be retained.

The executable pipeline uses \texttt{shap.GradientExplainer} on a thin
wrapper that exposes the breed-logit output of the VMGP model. The
explainer background is a deterministic stratified sample of
$n_{\mathrm{bg}}=100$ training individuals, and attributions are
computed on a second deterministic stratified sample of
$n_{\mathrm{eval}}=200$ training individuals (or on all training rows
when fewer are available). The sampling seed is fixed to $42$ and the
explainer seed is set to the same value. Attributions are taken in
absolute value and averaged across classes and evaluation samples; no
validation or test individual contributes to the background, the
evaluation subset or the aggregation.

\section{Explaining a self-supervised contrastive space}
\label{sec:contrastive_marker_attribution}

Attributing the contrastive model is conceptually more difficult. The
representation learner has no breed-classification head, so there is
no native breed logit to explain.

Integrated Gradients \cite{sundararajan2017axiomatic} remains
applicable to differentiable scalar functions of the model, but a
meaningful scalar target must first be defined from the learned
representation. This target may be related to a downstream
discriminative function or to a geometric property of the embedding,
depending on the final attribution design.

This distinction is important enough to make explicit. Integrated
Gradients is an attribution \emph{mechanism}; it does not by itself
determine what aspect of the contrastive representation should count
as breed relevance. The chosen target therefore becomes part of the
scientific method.

In the definitive protocol the scalar target is the cosine similarity
between the $L_2$-normalized embedding and the breed centroid of the
sample's own breed, where the centroid is the $L_2$-normalized mean of
the training embeddings of that breed. The baseline is the zero vector
in the four-channel one-hot genotype space, and the path integral is
approximated with $50$ steps. Per-marker per-sample attributions are
obtained as the $L_1$ sum over the four channels and aggregated as
described above. The supervised probe path implemented in the earlier
notebook is not used for the definitive RQ4 rankings.

\section{Population-genetics reference: fixation index}
\label{sec:fst_ranking}

A classical comparison is provided by locus-wise population
differentiation. For a SNP, $F_{ST}$ measures how strongly allele
frequencies differ among populations
\cite{wright1949genetical, weir1984estimating}.

Using $F_{ST}$ as a ranking has two useful properties. First, the score
has a direct population-genetics interpretation. Second, it does not
depend on the internal behavior of a neural network.

At the same time, a locus-wise statistic and a deep attribution method
should not be expected to produce identical rankings. The former
measures population differentiation at each marker, while the latter
describes relevance under a learned multivariate prediction or
representation. Their comparison is therefore informative precisely
because they are based on different principles.

\section{Machine-learning reference: Random Forest}
\label{sec:rf_ranking}

Random Forest provides a second reference \cite{breiman2001random}.
Unlike $F_{ST}$, it is explicitly predictive and can model non-linear
interactions between features.

A marker ranking can be obtained from a fitted Random Forest using a
chosen importance estimator. However, feature importance in tree
ensembles can be affected by correlated predictors and by the
estimator used \cite{strobl2007bias, nicodemus2010correlation}. The
final thesis should therefore report the exact importance definition
rather than referring generically to ``Random Forest importance''.

The frozen Random Forest ranking uses impurity-based (Gini) importance
from $200$ trees with a fixed random seed ($42$), fitted only on the
training partition of each fold. Because impurity importance is known
to favour markers with high minor-allele frequency and correlated
predictors, the thesis treats this ranking as a predictive
machine-learning reference rather than as an unbiased estimate of
biological relevance.

\section{Constructing nested reduced panels}
\label{sec:reduced_panel_construction}

Once a ranking
$r=(j_1,j_2,\ldots,j_P)$
has been produced, reduced marker panels can be constructed
progressively,

\begin{equation}
    \mathcal{P}_K
    =
    \{j_1,j_2,\ldots,j_K\}.
\end{equation}

The nested construction is useful because every larger panel contains
the markers used by the smaller one. Changes in downstream performance
can therefore be interpreted as the effect of adding progressively
lower-ranked markers rather than as the difference between unrelated
subsets.

The full post-QC panel remains the reference point. The purpose of RQ3
is not to maximize performance at any cost, but to identify how
aggressively the marker set can be reduced while retaining a
pre-specified fraction of the full-panel classification performance.

The definitive grid is $K \in \{50, 100, 200, 500, 1000, 2000, 5000\}$,
fixed before inspecting any downstream result. The primary retention
tolerance is $\epsilon_P = 0.02$, with $\epsilon_P = 0.05$ reported as
a secondary reference. For every ranking method, the smallest grid
value satisfying the criterion is reported as $P_{\min}$.

\section{Leakage prevention during marker selection}
\label{sec:marker_selection_leakage}

Marker selection is itself a learning procedure. If a held-out
individual's breed label or genotype pattern influences the ranking,
downstream validation is no longer independent.

For this reason, every ranking considered in the thesis must be fitted
or constructed inside the relevant training partition. The validation
partition is used only to evaluate the reduced panel that results from
that ranking. The held-out test set remains unavailable throughout
marker-ranking development.

This requirement applies equally to apparently ``classical'' methods
and neural attributions. $F_{ST}$, Random Forest importance, SHAP,
Integrated Gradients, and any derived aggregation must obey the same
information boundary.

\section{Ranking stability and downstream utility}
\label{sec:ranking_stability_utility}

Two complementary questions can be asked about a marker ranking.

The first is \emph{stability}: do different cross-validation folds or
sampling replicates tend to identify similar high-ranked markers? This
can be investigated using overlap measures or rank correlations.

The second is \emph{utility}: does a panel built from the ranking
preserve breed-classification performance when evaluated on unseen
individuals?

Stability without predictive utility is not sufficient, because a
ranking could be reproducible but uninformative. Conversely, high
predictive utility with unstable marker identities may indicate that
several correlated genomic subsets contain interchangeable
information.

Considering both views is especially important in genomic data, where
LD can make multiple marker panels functionally similar even when
their exact SNP identities differ.

\section{What this chapter does not claim}
\label{sec:attribution_not_causality}

The marker-ranking analysis is computational. A high attribution,
$F_{ST}$ value, or Random Forest importance score does not by itself
establish that a locus is causal for breed identity, nor that it
should be interpreted as a functional variant.

The output of the framework is instead a ranked set of candidate
markers whose value is evaluated through reproducibility and
downstream discrimination. Biological interpretation of individual
loci would require additional genomic annotation and domain-specific
validation beyond the scope of the present computational study.
